\section{Circuit Complexity --- Parity}
%%% generated by the blueprint pipeline from TCSlib.Complexity.CircuitComplexity.Parity
\subsection*{Circuit Complexity --- Parity: Introduction}

Declarations of the Lean module \texttt{TCSlib.Complexity.CircuitComplexity.Parity}.

\subsection*{Circuit Complexity --- Parity: Results}

\begin{definition}[XOR circuit construction]
\lean{BoolCircuit.xorNode}
\label{BoolCircuit.xorNode}
\leanok
\uses{BoolCircuit.TreeCircuit}
For a natural number $n$ and pairs $p=(p_1,p_2)$ and $q=(q_1,q_2)$ of Boolean circuit
trees on $n$ variables, define a pair of circuit trees whose first component is
$(p_1\land q_2)\lor(p_2\land q_1)$ and whose second component is $(p_1\land
q_1)\lor(p_2\land q_2)$, with each conjunction and disjunction represented by the
corresponding circuit node.
\end{definition}

\begin{definition}[Dual Boolean circuits]
\lean{BoolCircuit.IsDual}
\statementsource{Analysis-of-Boolean-Functions-by-Ryan-ODonnell}{pdf-ffbe9c7b3306-p094-b005}
\label{BoolCircuit.IsDual}
\leanok
\uses{BoolCircuit.TreeCircuit, BoolCircuit.TreeCircuit.eval}
For a function $x$ assigning a Boolean value to each of $n$ variables and a pair $(C,D)$
of Boolean circuit trees on those variables, the circuits are dual at $x$ when $D(x)=
eg C(x)$.
\end{definition}

\begin{theorem}[Evaluation of one XOR step]
\lean{BoolCircuit.xorNode_eval}
\label{BoolCircuit.xorNode_eval}
\leanok
\uses{BoolCircuit.IsDual, BoolCircuit.TreeCircuit, BoolCircuit.TreeCircuit.eval, BoolCircuit.xorNode}
\difficulty{4}
Let $n$ be a nonnegative integer, let $x$ assign Boolean values to the $n$ variables,
and let $p,q$ be pairs of Boolean circuit trees on those variables. Suppose the second
circuit in each pair evaluates under $x$ to the negation of the first. Then the first
circuit in the pair produced by one XOR step evaluates to the exclusive OR of the first
circuits in $p$ and $q$.
\end{theorem}

\begin{theorem}[Complementarity after one XOR step]
\lean{BoolCircuit.xorNode_isDual}
\label{BoolCircuit.xorNode_isDual}
\leanok
\uses{BoolCircuit.IsDual, BoolCircuit.TreeCircuit, BoolCircuit.TreeCircuit.eval, BoolCircuit.xorNode}
\difficulty{4}
Let $n$ be a nonnegative integer, let $x$ assign Boolean values to the $n$ variables,
and let $p,q$ be pairs of Boolean circuit trees on those variables. If the two circuits
in each pair evaluate under $x$ to complementary values, then the two circuits in the
pair produced by one XOR step also evaluate under $x$ to complementary values.
\end{theorem}

\begin{definition}[Parity of circuit values]
\lean{BoolCircuit.xorAll}
\label{BoolCircuit.xorAll}
\leanok
\uses{BoolCircuit.TreeCircuit, BoolCircuit.TreeCircuit.eval}
For a Boolean assignment $x$ to the $n$ variables and a finite list of pairs $(C_i,D_i)$
of Boolean circuit trees, define the result to be the exclusive OR of the values of the
first circuits $C_i$, evaluated at $x$, with the empty list giving false.
\end{definition}

\begin{definition}[Maximum depth of circuit pairs]
\lean{BoolCircuit.pairDepth}
\label{BoolCircuit.pairDepth}
\leanok
\uses{BoolCircuit.TreeCircuit, BoolCircuit.TreeCircuit.depth}
For a finite list of pairs $(C_i,D_i)$ of Boolean circuit trees, define the pair depth
to be the maximum of $0$ and the depths of all circuits occurring in the pairs.
\end{definition}

\begin{definition}[Maximum fan-in of circuit pairs]
\lean{BoolCircuit.pairFanin}
\label{BoolCircuit.pairFanin}
\leanok
\uses{BoolCircuit.TreeCircuit, BoolCircuit.TreeCircuit.maxFanin}
For a finite list of pairs $(C_i,D_i)$ of Boolean circuit trees, define the pair fan-in
to be the maximum of $0$ and the maximum fan-ins of all circuits occurring in the pairs.
\end{definition}

\begin{theorem}[Xor over a list with a leading pair]
\lean{BoolCircuit.xorAll_cons}
\label{BoolCircuit.xorAll_cons}
\leanok
\uses{BoolCircuit.TreeCircuit, BoolCircuit.TreeCircuit.eval, BoolCircuit.xorAll}
\difficulty{1}
Let $n$ be a nonnegative integer, let $x$ assign Boolean values to its variables, let
$p$ be a pair of Boolean circuit trees on those variables, and let $P$ be a finite list
of such pairs. The exclusive-or of the first circuit in each pair, accumulated over the
list $p::P$ with false as the initial value, is the exclusive-or of the first circuit in
$p$ with the corresponding accumulation over $P$.
\end{theorem}

\begin{theorem}[Pair depth of a list with a leading pair]
\lean{BoolCircuit.pairDepth_cons}
\label{BoolCircuit.pairDepth_cons}
\leanok
\uses{BoolCircuit.TreeCircuit, BoolCircuit.TreeCircuit.depth, BoolCircuit.pairDepth}
\difficulty{1}
Let $p$ be a pair of Boolean circuits and let $ps$ be a list of such pairs. The pair
depth of the list beginning with $p$ is the maximum of the depths of the two circuits in
$p$ and the pair depth of $ps$.
\end{theorem}

\begin{theorem}[Pair fanin of a list with a head]
\lean{BoolCircuit.pairFanin_cons}
\label{BoolCircuit.pairFanin_cons}
\leanok
\uses{BoolCircuit.TreeCircuit, BoolCircuit.TreeCircuit.maxFanin, BoolCircuit.pairFanin}
\difficulty{1}
For every natural number $n$, every pair of circuits $p$ on $n$ variables, and every
finite list $ps$ of such pairs, the pair fanin of $p$ followed by $ps$ is the maximum of
the maximum fanin of the two circuits in $p$ and the pair fanin of $ps$.
\end{theorem}

\begin{theorem}[Depth of the XOR pair construction]
\lean{BoolCircuit.pairDepth_xorNode}
\label{BoolCircuit.pairDepth_xorNode}
\leanok
\uses{BoolCircuit.TreeCircuit, BoolCircuit.TreeCircuit.depth, BoolCircuit.TreeCircuit.depth_node, BoolCircuit.TreeCircuit.maxDepth_cons, BoolCircuit.TreeCircuit.maxDepth_nil, BoolCircuit.xorNode}
\difficulty{3}
For every natural number $n$ and any two pairs of circuits $(p_1,p_2)$ and $(q_1,q_2)$
on $n$ variables, the maximum depth of the two circuits produced by the XOR pair
construction is at most two plus the maximum depth among the four input circuits.
\end{theorem}

\begin{theorem}[Fanin bound for the XOR pair construction]
\lean{BoolCircuit.pairFanin_xorNode}
\label{BoolCircuit.pairFanin_xorNode}
\leanok
\uses{BoolCircuit.TreeCircuit, BoolCircuit.TreeCircuit.maxFanin, BoolCircuit.TreeCircuit.maxFaninL_cons, BoolCircuit.TreeCircuit.maxFaninL_nil, BoolCircuit.TreeCircuit.maxFanin_node, BoolCircuit.xorNode}
\difficulty{3}
For every natural number $n$ and any two pairs of circuits $(p_1,p_2)$ and $(q_1,q_2)$
on $n$ variables, the maximum fanin of the two circuits produced by the XOR pair
construction is at most the maximum of $2$ and the maximum fanin among the four input
circuits.
\end{theorem}

\begin{definition}[Pairwise exclusive-OR construction]
\lean{BoolCircuit.xorPairUp}
\label{BoolCircuit.xorPairUp}
\leanok
\uses{BoolCircuit.TreeCircuit, BoolCircuit.xorNode}
Given pairs of Boolean circuit trees $(A,B)$ and $(C,D)$, define their combined pair to
be $((A\land D)\lor(B\land C),\,(A\land C)\lor(B\land D))$, where each conjunction and
disjunction is represented by a circuit node.
\end{definition}

\begin{theorem}[Length after pairing circuit pairs]
\lean{BoolCircuit.length_xorPairUp}
\label{BoolCircuit.length_xorPairUp}
\leanok
\uses{BoolCircuit.TreeCircuit, BoolCircuit.xorPairUp}
\difficulty{4}
Let $n\in\mathbb{N}$ and let the entries of $ps$ be pairs of Boolean circuit trees on
$n$ variables. Pairing consecutive entries by the specified operation, while retaining a
final unpaired entry if present, produces a list of length
$\left\lfloor\frac{|ps|+1}{2}\right\rfloor$.
\end{theorem}

\begin{theorem}[Duality Preserved by Pairwise XOR Construction]
\lean{BoolCircuit.isDual_xorPairUp}
\label{BoolCircuit.isDual_xorPairUp}
\leanok
\uses{BoolCircuit.IsDual, BoolCircuit.TreeCircuit, BoolCircuit.xorNode_isDual, BoolCircuit.xorPairUp}
\difficulty{5}
Let $n$ be a natural number, fix an assignment of Boolean values to its $n$ variables,
and let a pair of circuits be called dual when the second evaluates to the Boolean
complement of the first. Given any list of circuit pairs all of whose pairs are dual,
combine consecutive pairs $(p,q)$ into the pair whose first circuit is $(p_1\land
q_2)\lor(p_2\land q_1)$ and whose second circuit is $(p_1\land q_1)\lor(p_2\land q_2)$;
an unpaired final circuit pair is left as it is. Every pair in the resulting list is
dual.
\end{theorem}

\begin{theorem}[Pairing preserves accumulated xor]
\lean{BoolCircuit.xorAll_xorPairUp}
\label{BoolCircuit.xorAll_xorPairUp}
\leanok
\uses{BoolCircuit.IsDual, BoolCircuit.TreeCircuit, BoolCircuit.xorAll, BoolCircuit.xorAll_cons, BoolCircuit.xorNode_eval, BoolCircuit.xorPairUp}
\difficulty{4}
Let $n$ be a nonnegative integer, let $x$ assign Boolean values to its variables, and
let $P$ be a finite list of pairs of Boolean circuit trees on those variables. Suppose
that in every pair the second circuit evaluates to the negation of the first under $x$.
Replacing successive pairs in $P$ by the paired circuit construction preserves the
exclusive-or of the first circuit in each pair, accumulated with false as the initial
value.
\end{theorem}

\begin{theorem}[Depth bound for pairwise XOR construction]
\lean{BoolCircuit.pairDepth_xorPairUp}
\label{BoolCircuit.pairDepth_xorPairUp}
\leanok
\uses{BoolCircuit.TreeCircuit, BoolCircuit.pairDepth, BoolCircuit.pairDepth_cons, BoolCircuit.pairDepth_xorNode, BoolCircuit.xorPairUp}
\difficulty{4}
For every natural number $n$ and every finite list of pairs of circuits on $n$
variables, applying the pairwise XOR construction yields a list whose pair depth is at
most two plus the original list's pair depth.
\end{theorem}

\begin{theorem}[Fanin bound for pairwise XOR construction]
\lean{BoolCircuit.pairFanin_xorPairUp}
\label{BoolCircuit.pairFanin_xorPairUp}
\leanok
\uses{BoolCircuit.TreeCircuit, BoolCircuit.pairFanin, BoolCircuit.pairFanin_cons, BoolCircuit.pairFanin_xorNode, BoolCircuit.xorPairUp}
\difficulty{4}
For every natural number $n$ and every finite list of pairs of circuits on $n$
variables, applying the pairwise XOR construction yields a list whose pair fanin is at
most the maximum of $2$ and the original list's pair fanin.
\end{theorem}

\begin{definition}[Fuel-bounded exclusive-OR construction]
\lean{BoolCircuit.xorFuel}
\label{BoolCircuit.xorFuel}
\leanok
\uses{BoolCircuit.TreeCircuit, BoolCircuit.xorPairUp}
For a natural-number fuel and a finite list of pairs of Boolean circuit trees, define a
pair recursively as follows. With zero fuel, or with an empty list, return the pair
consisting of the empty disjunction circuit and the empty conjunction circuit; with
positive fuel and a singleton list, return its sole pair; with positive fuel and at
least two pairs, combine the first two by the pairwise construction and recurse on the
resulting list with one less fuel.
\end{definition}

\begin{theorem}[Evaluation of the XOR reduction]
\lean{BoolCircuit.xorFuel_eval}
\label{BoolCircuit.xorFuel_eval}
\leanok
\uses{BoolCircuit.IsDual, BoolCircuit.TreeCircuit, BoolCircuit.TreeCircuit.eval, BoolCircuit.isDual_xorPairUp, BoolCircuit.length_xorPairUp, BoolCircuit.xorAll, BoolCircuit.xorAll_xorPairUp, BoolCircuit.xorFuel, BoolCircuit.xorPairUp}
\difficulty{4}
Let $n,k$ be nonnegative integers, let $x$ be an assignment of Boolean values to the $n$
variables, and let $ps$ be a list of pairs of Boolean circuit trees on those variables
with $|ps|\le k$. Suppose that in every pair the second circuit evaluates under $x$ to
the negation of the first. Then the first circuit returned by the fuel-bounded XOR
reduction evaluates to the iterated exclusive OR of the first circuits in $ps$, and the
returned pair consists of circuits whose values under $x$ are complements.
\end{theorem}

\begin{theorem}[Depth bound for XOR reduction]
\lean{BoolCircuit.xorFuel_depth}
\label{BoolCircuit.xorFuel_depth}
\leanok
\uses{BoolCircuit.TreeCircuit, BoolCircuit.TreeCircuit.depth, BoolCircuit.TreeCircuit.depth_node, BoolCircuit.TreeCircuit.maxDepth_nil, BoolCircuit.length_xorPairUp, BoolCircuit.pairDepth, BoolCircuit.pairDepth_cons, BoolCircuit.pairDepth_xorPairUp, BoolCircuit.xorFuel, BoolCircuit.xorPairUp}
\difficulty{5}
Let $n,k$ be nonnegative integers and let $ps$ be a list of pairs of Boolean circuit
trees on $n$ variables, with $|ps|\le k$. Define the pair depth of $ps$ as the maximum,
over its pairs, of the larger depth in each pair, taking this maximum to be $0$ for the
empty list. The maximum depth of the two circuits returned by the fuel-bounded XOR
reduction is at most the pair depth of $ps$ plus $2\lceil\log_2 |ps|\rceil+2$, where the
ceiling is understood as the least integer $m$ such that $|ps|\le 2^m$.
\end{theorem}

\begin{theorem}[Fanin bound for XOR reduction]
\lean{BoolCircuit.xorFuel_fanin}
\label{BoolCircuit.xorFuel_fanin}
\leanok
\uses{BoolCircuit.TreeCircuit, BoolCircuit.TreeCircuit.maxFanin, BoolCircuit.TreeCircuit.maxFaninL_nil, BoolCircuit.TreeCircuit.maxFanin_node, BoolCircuit.length_xorPairUp, BoolCircuit.pairFanin, BoolCircuit.pairFanin_cons, BoolCircuit.pairFanin_xorPairUp, BoolCircuit.xorFuel, BoolCircuit.xorPairUp}
\difficulty{4}
Let $n,k$ be nonnegative integers and let $ps$ be a list of pairs of Boolean circuit
trees on $n$ variables, with $|ps|\le k$. Define the pair fanin of $ps$ as the maximum
fanin among all circuits in its pairs, taking it to be $0$ for the empty list. The
maximum fanin of the two circuits returned by the fuel-bounded XOR reduction is at most
the larger of $2$ and the pair fanin of $ps$.
\end{theorem}

\begin{definition}[Balanced exclusive-OR construction]
\lean{BoolCircuit.xorTree}
\label{BoolCircuit.xorTree}
\leanok
\uses{BoolCircuit.TreeCircuit, BoolCircuit.xorFuel}
For a finite list of pairs of Boolean circuit trees, define its resulting pair by
applying the fuel-bounded construction with fuel equal to the length of the list.
\end{definition}

\begin{definition}[Positive and negative literal pairs]
\lean{BoolCircuit.parityPairs}
\label{BoolCircuit.parityPairs}
\leanok
\uses{BoolCircuit.TreeCircuit}
For each natural number $m$, define a list of $m$ pairs of literal circuits, one pair
for each variable index $i<m$: the pair consists of the positive literal for variable
$i$ and its negative literal.
\end{definition}

\begin{theorem}[Duality of Literal Pairs]
\lean{BoolCircuit.isDual_parityPairs}
\label{BoolCircuit.isDual_parityPairs}
\leanok
\uses{BoolCircuit.IsDual, BoolCircuit.Lit.eval, BoolCircuit.TreeCircuit.eval, BoolCircuit.parityPairs}
\difficulty{4}
Let $n$ be a natural number and assign a Boolean value to each of its $n$ variables. For
each variable, form the pair consisting of its positive literal and its negated literal.
In every such pair, the value of the second circuit is the Boolean complement of the
value of the first.
\end{theorem}

\begin{theorem}[Depth of the parity literal pairs]
\lean{BoolCircuit.pairDepth_parityPairs}
\label{BoolCircuit.pairDepth_parityPairs}
\leanok
\uses{BoolCircuit.TreeCircuit.depth, BoolCircuit.pairDepth, BoolCircuit.pairDepth_cons, BoolCircuit.parityPairs}
\difficulty{4}
For every natural number $n$, the list of $n$ pairs of literals, with the two opposite
polarities at each variable index, has pair depth $0$.
\end{theorem}

\begin{theorem}[Fanin of the parity literal pairs]
\lean{BoolCircuit.pairFanin_parityPairs}
\label{BoolCircuit.pairFanin_parityPairs}
\leanok
\uses{BoolCircuit.TreeCircuit.maxFanin, BoolCircuit.pairFanin, BoolCircuit.pairFanin_cons, BoolCircuit.parityPairs}
\difficulty{4}
For every natural number $n$, the list of $n$ pairs of literals, with the two opposite
polarities at each variable index, has pair fanin $0$.
\end{theorem}

\begin{theorem}[Number of parity pairs]
\lean{BoolCircuit.length_parityPairs}
\label{BoolCircuit.length_parityPairs}
\leanok
\uses{BoolCircuit.parityPairs}
\difficulty{1}
For every natural number $n$, the list of pairs consisting, for each of the $n$
variables, of its positive and negative literals has length $n$.
\end{theorem}

\begin{theorem}[Parity pairs evaluate to input parity]
\lean{BoolCircuit.xorAll_parityPairs}
\label{BoolCircuit.xorAll_parityPairs}
\leanok
\uses{BoolCircuit.Lit.eval, BoolCircuit.TreeCircuit.eval, BoolCircuit.parityPairs, BoolCircuit.xorAll, BoolCircuit.xorAll_cons}
\difficulty{4}
Let $n$ be a nonnegative integer and let $x$ assign Boolean values to its variables. The
exclusive-or accumulated over the list of pairs consisting, for each variable, of its
positive and negative literal evaluates to the exclusive-or of the input values,
accumulated over all variable indices with false as the initial value.
\end{theorem}

\begin{definition}[Parity circuit]
\lean{BoolCircuit.parityCircuit}
\label{BoolCircuit.parityCircuit}
\leanok
\uses{BoolCircuit.TreeCircuit, BoolCircuit.parityPairs, BoolCircuit.xorTree}
For each natural number $m$, define a Boolean circuit tree to be the first circuit in
the balanced exclusive-OR construction applied to the list of positive and negative
literal pairs for the $m$ variables.
\end{definition}

\begin{theorem}[Evaluation of the parity circuit]
\lean{BoolCircuit.parityCircuit_eval}
\label{BoolCircuit.parityCircuit_eval}
\leanok
\uses{BoolCircuit.TreeCircuit.eval, BoolCircuit.isDual_parityPairs, BoolCircuit.parityCircuit, BoolCircuit.parityPairs, BoolCircuit.xorAll_parityPairs, BoolCircuit.xorFuel_eval, BoolCircuit.xorTree}
\difficulty{3}
For every natural number $n$ and every assignment $x:\text{Fin}(n)\to\text{Bool}$, the
parity circuit evaluates to the iterated exclusive-or of $x(0),\ldots,x(n-1)$, with
initial value false; equivalently, its value is the parity of the assigned bits.
\end{theorem}

\begin{theorem}[Fanin bound for the parity circuit]
\lean{BoolCircuit.parityCircuit_maxFanin_le}
\label{BoolCircuit.parityCircuit_maxFanin_le}
\leanok
\uses{BoolCircuit.TreeCircuit.maxFanin, BoolCircuit.pairFanin_parityPairs, BoolCircuit.parityCircuit, BoolCircuit.parityPairs, BoolCircuit.xorFuel_fanin}
\difficulty{3}
For every natural number $n$, the maximum fanin of the parity circuit on $n$ Boolean
variables is at most $2$.
\end{theorem}

\begin{theorem}[Depth bound for the parity circuit]
\lean{BoolCircuit.parityCircuit_depth_le}
\label{BoolCircuit.parityCircuit_depth_le}
\leanok
\uses{BoolCircuit.TreeCircuit.depth, BoolCircuit.length_parityPairs, BoolCircuit.pairDepth, BoolCircuit.pairDepth_parityPairs, BoolCircuit.parityCircuit, BoolCircuit.parityPairs, BoolCircuit.xorFuel_depth}
\difficulty{4}
For every natural number $n$, the depth of the parity circuit on $n$ Boolean variables
is at most $4(\lfloor\log_2 n\rfloor+1)$, with the logarithm interpreted as the
natural-number logarithm used in the bound.
\end{theorem}

\begin{theorem}[Size bound for the parity circuit]
\lean{BoolCircuit.parityCircuit_size_le}
\label{BoolCircuit.parityCircuit_size_le}
\leanok
\uses{BoolCircuit.TreeCircuit.depth, BoolCircuit.TreeCircuit.size, BoolCircuit.TreeCircuit.size_succ_le_two_pow, BoolCircuit.parityCircuit, BoolCircuit.parityCircuit_depth_le, BoolCircuit.parityCircuit_maxFanin_le}
\difficulty{5}
For every natural number $n$, the size of the parity circuit on $n$ Boolean variables is
at most $32(n+1)^4$.
\end{theorem}

\begin{theorem}[Value of the zero-variable parity circuit]
\lean{BoolCircuit.parityCircuit_eval_zero}
\label{BoolCircuit.parityCircuit_eval_zero}
\leanok
\uses{BoolCircuit.TreeCircuit.eval, BoolCircuit.parityCircuit, BoolCircuit.parityCircuit_eval}
\difficulty{3}
Let $x$ be an assignment of Boolean values to the zero variables. The parity circuit on
zero variables evaluates to false under $x$.
\end{theorem}

\begin{theorem}[Evaluation of the one-variable parity circuit]
\lean{BoolCircuit.parityCircuit_eval_one}
\label{BoolCircuit.parityCircuit_eval_one}
\leanok
\uses{BoolCircuit.TreeCircuit.eval, BoolCircuit.parityCircuit, BoolCircuit.parityCircuit_eval}
\difficulty{3}
For every Boolean assignment $x : \{0\} \to \{\mathrm{false},\mathrm{true}\}$ to the
single variable, the parity circuit on one variable evaluates to $x(0)$.
\end{theorem}

\begin{definition}[Odd-parity Boolean words]
\lean{Language.parity}
\statementsource{Vad12}{pdf-bf2ce5422088-p240-b002}
\statementsource{arora-barak-ch13-circuit-lowerbounds}{pdf-52929dbaa962-p001-b004}
\statementsource{Analysis-of-Boolean-Functions-by-Ryan-ODonnell}{
  pdf-ffbe9c7b3306-p023-b004,
  pdf-ffbe9c7b3306-p148-b004
}
\statementsource{goldreich-approximation-method}{pdf-988dc9f0ae91-p004-b004}
\label{Language.parity}
\leanok
The language consists of all finite Boolean words containing an odd number of true
entries.
\end{definition}

\begin{theorem}[Mem parity iff]
\lean{Language.mem_parity_iff}
\statementsource{goldreich-approximation-method}{pdf-988dc9f0ae91-p004-b004}
\label{Language.mem_parity_iff}
\leanok
\uses{Language.parity}
\difficulty{5}
% dataset-pipeline: stub -- prose pending build_informal_statements.py
Statement pending regeneration.
\end{theorem}

\begin{theorem}[Parity lies in tree log-depth circuit complexity]
\lean{Language.parity_inTreeNC_one}
\label{Language.parity_inTreeNC_one}
\leanok
\uses{BoolCircuit.TreeCircuitFamily.mem_language_iff, BoolCircuit.parityCircuit, BoolCircuit.parityCircuit_depth_le, BoolCircuit.parityCircuit_eval, BoolCircuit.parityCircuit_maxFanin_le, BoolCircuit.parityCircuit_size_le, Language.InTreeNC, Language.mem_parity_iff, Language.parity}
\difficulty{3}
The language of Boolean lists containing an odd number of true entries is recognized by
a family of Boolean formula trees with fan-in at most two, size bounded by a polynomial
in the input length, and depth bounded by a constant multiple of $\log_2 n+1$. In other
words, parity belongs to the class defined by these conditions with depth exponent $1$.
\end{theorem}

\begin{theorem}[Parity lies in log-depth circuit complexity]
\lean{Language.parity_inNC_one}
\label{Language.parity_inNC_one}
\leanok
\uses{Language.InNC, Language.InTreeNC.inNC, Language.parity, Language.parity_inTreeNC_one}
\difficulty{1}
The language of Boolean lists containing an odd number of true entries is recognized by
a family of acyclic Boolean circuits with fan-in at most two, size bounded by a
polynomial in the input length, and depth bounded by a constant multiple of $\log_2
n+1$. In other words, parity belongs to the class defined by these conditions with depth
exponent $1$.
\end{theorem}
